% =====================================================================
% Chapter 1: Introduction
% =====================================================================
\chapter{Introduction}
\label{ch:introduction}

\section{Overview}
\label{sec:overview}

Football (soccer) is the most widely followed sport in the world, and
the broadcast analysis of matches generates a substantial volume of
visual data each season. Identifying individual players in this data
stream is a fundamental requirement for many downstream applications,
including automated highlight generation, tactical analysis, player
performance monitoring, and real-time statistics overlays for
television audiences. Although player identity can, in principle, be
inferred from facial features, gait patterns, or team-roster information,
the jersey number printed on a player's shirt remains the most
direct, universally interpretable, and camera-accessible cue for
player identification in broadcast footage.

Despite appearing straightforward at first glance, automatic jersey
number recognition is a genuinely difficult computer vision task. The
digit regions occupy only a small fraction of a typical video frame,
are routinely blurred by player motion, may be partly hidden by the
ball, other players, or a player's own limbs, and can appear at
steep viewing angles due to broadcast camera placement. Jersey color,
fabric texture, typeface variation across teams, and the extremely low
resolution of digit patches in standard broadcast video complicate the
task further. Simply passing a full frame or an uncropped player
bounding box to a general-purpose OCR engine consistently yields poor
recognition results.

This thesis addresses these challenges by designing, implementing, and
evaluating a complete jersey number recognition system that functions in
two complementary modes. The first is image-based recognition,
evaluated on the SoccerNet SN-Jersey dataset, which provides cropped
player images grouped by player identity together with ground-truth
jersey numbers. The second is video-based recognition, evaluated on ten
broadcast football clips from the SoccerNet tracking dataset (videos
sourced from the Hugging Face SoccerNet SN-PCBAS release and from
Rutracker.org), where players must first be detected and tracked across
frames before OCR evidence can be accumulated over time. Both pipelines
share a conceptual foundation --- player or torso localization followed
by image enhancement and OCR --- but differ substantially in
implementation due to the different nature of their inputs: isolated
player crops versus continuous video streams.

\section{Motivation}
\label{sec:motivation}

Several practical and academic considerations motivated this research.
From an applied standpoint, automated player identification is a
fundamental building block for modern broadcast and club analytics
platforms, since manual annotation of jersey numbers across hours of
footage is both prohibitively expensive and unacceptably slow for
operational use. General-purpose OCR engines, including Tesseract and
out-of-the-box EasyOCR, are designed for scanned documents and perform
poorly on the small, distorted, stylized digits encountered on curved
fabric under variable stadium lighting, so a domain-specific pipeline
with targeted pre-processing and a multi-variant recognition strategy
is essential to obtain acceptable accuracy. Precise spatial
localization is equally important: jersey digits typically occupy only
a tiny region of a broadcast frame, and without accurate player
detection and torso localization, OCR cannot be reliably directed to
the correct region of interest. In video, temporal redundancy offers a
further opportunity, since a single frame often yields no usable OCR
result due to motion blur or occlusion, while a given player typically
appears across many consecutive frames; exploiting this temporal
information through a voting mechanism substantially improves overall
recognition reliability. Finally, jersey number recognition integrates
several core deep-learning and computer vision topics --- object
detection, pose estimation, multi-object tracking, image enhancement,
OCR, and binary classification --- making it an ideal capstone project
for demonstrating practical, integrated skills.

\section{Problem Statement}
\label{sec:problem-statement}

Given either a static image of a football player or a broadcast video
clip containing multiple players, the goal of this thesis is to
automatically detect and localize each player, tracking each one
persistently across frames using a unique identity when the input is
video; derive the torso region of each player, which is the most
likely location of the jersey number; enhance the localized region to
compensate for blur, poor contrast, and noise; recognize the one- or
two-digit jersey number from the enhanced torso crop using OCR; and,
for video, aggregate per-frame OCR predictions into a single, reliable
jersey-number estimate per tracked player through a temporal voting
mechanism. The system must additionally provide a means of
distinguishing legible jersey crops from illegible ones, so that
downstream components receive information about prediction confidence.

\section{Objectives}
\label{sec:objectives}

The specific objectives of this thesis are, first, to design and
implement a YOLOv8-based player detection and YOLOv8x-Pose-based torso
localization pipeline for the image-based SoccerNet SN-Jersey dataset;
second, to build a global image-enhancement module combining
denoising, CLAHE, bilateral filtering, and unsharp masking to improve
jersey patch quality prior to OCR; third, to develop a multi-variant
OCR pipeline using EasyOCR that incorporates dynamic zoom, multiple
binarization strategies, per-digit isolation, projection-based digit
splitting, rotation correction, and shape-based digit disambiguation
for the digit pairs 0/6/8/9 and 1/7; fourth, to train a ResNet-18-based
binary legibility classifier that identifies whether a cropped jersey
region is readable, using weak labels derived automatically from OCR
confidence scores; fifth, to design and implement a video-based
pipeline that uses a specialist football detection model (YOLOv5) with
ByteTrack multi-object tracking, a temporal role voter for referee
identification, and a PARSeq-based jersey OCR processor with automatic
fallback to EasyOCR; sixth, to implement a temporal jersey voting
mechanism that aggregates per-frame jersey number predictions over a
sliding window per track, and to evaluate it against manually
annotated ground truth across ten broadcast clips; and finally, to
evaluate both pipelines using quantitative metrics including confusion
matrices, player-level accuracy, tracking metrics (MOTA / IDF1), OCR
success rates before and after temporal voting, and processing-time
breakdowns.

\section{Scope of the Thesis}
\label{sec:scope}

This thesis focuses specifically on the visual recognition of jersey
numbers and does not attempt to solve the full player re-identification
problem. The image-based pipeline is evaluated on a curated subset of
353 player-image folders from the SoccerNet SN-Jersey-2023 dataset (269
training, 84 test). The video-based pipeline is evaluated on ten
broadcast clips (001--010.mp4), each 900 frames in length, collected
from the Hugging Face SoccerNet SN-PCBAS dataset (Videos 001--007) and
Rutracker.org (Videos 008--010), and evaluated against manually
annotated ground-truth CSV files. The system assumes jersey numbers
consist of one or two digits in the range 1--99. Multi-camera fusion,
three-digit numbers, and non-numeric jersey markings are outside the
scope of this work.

\section{Organization of the Report}
\label{sec:organization}

Chapter~\ref{ch:literature} reviews related work in jersey number
recognition, scene text recognition, and sports analytics.
Chapter~\ref{ch:methodology} describes the datasets, pre-processing
pipeline, model architectures, algorithmic formulation, and evaluation
metrics. Chapter~\ref{ch:system-design} presents the implementation
details for both pipelines, including workflow diagrams.
Chapter~\ref{ch:results} presents experimental results, with a
comparative discussion against prior work from the literature.
Chapter~\ref{ch:conclusion} concludes the thesis and outlines
directions for future work.
